% ===========================================================================
%  板块四 · 求最值：一正二定三相等（卡片 4 + 题 8–13）
% ===========================================================================
\section{求最值：一正二定三相等}

\begin{strand}
\keyword{本板块主线}：用不等式求最值，成立与否取决于等号能否取到。
\end{strand}

\block{知识精讲：不等式为什么能求最值}

此前求最值有两条路：几何方法（如「将军饮马」，三点共线时最短）与函数方法（如二次函数配方）。基本不等式提供第三条路，其逻辑必须想清楚：

\begin{strand}
$y\ge 4$ 只说明 $y$ 不会小于 4，\textbf{不等于}说「$y$ 的最小值是 4」——$y=5$ 也满足 $y\ge 4$。只有当 $y=4$ 确实取得到（等号成立），4 才是最小值。
\end{strand}

因此用基本不等式求最值要做三项检查，\keyword{口诀：一正、二定、三相等}：

\begin{center}
\begin{tabularx}{\linewidth}{@{}l L@{}}
\toprule
\multicolumn{1}{l}{\color{indigo}检查} &
\multicolumn{1}{l}{\color{indigo}原因} \\
\midrule
\textbf{一正}：各项为正 & 基本不等式的前提 \\
\textbf{二定}：和或积为定值 & $x+y\ge 2\sqrt{xy}$ 的右边必须是常数才构成下界；积 $xy=P$ 为定值 $\Rightarrow$ 和的最小值 $2\sqrt{P}$；和 $x+y=S$ 为定值 $\Rightarrow$ 积的最大值 $\dfrac{S^2}{4}$ \\
\textbf{三相等}：等号取得到 & 取不到则该界不是最值（题 13 专门辨析） \\
\bottomrule
\end{tabularx}
\end{center}

与板块二的图对照：半弦 $DE=\sqrt{ab}$ 何时最长？垂足 $E$ 与圆心重合、$a=b$ 之时。\textbf{最值出现在等号成立的位置}，几何与代数一致。

% 注释开关：易错提醒（自总结版可 \iffalse … \fi 注释）
\textbf{易错提醒。}
\begin{itemize}
  \item 右边不是定值时，「$\ge$」不构成最小值结论。$x^2+\dfrac1x\ge 2\sqrt{x}$ 成立，但 $2\sqrt{x}$ 随 $x$ 变化，从中读不出最小值（题 13 D 即此错误）。
  \item 负变量先转正：$x<0$ 时对 $-x>0$ 使用基本不等式（题 9）。
  \item 取等的 $x$ 必须落在定义域内，代回检验是必要步骤。
\end{itemize}

\begin{problem}{题 8 （直接使用）}
当 $x>0$ 时，求 $y=4x+\dfrac{12}{x}$ 的最小值，并指出取到最小值时的 $x$。
\hint{两项之积是否为定值？计算 $4x\cdot\dfrac{12}{x}$。}
\end{problem}

\begin{problem}{题 9 （负变量先转正）}
已知 $x<0$，求 $y=x+\dfrac{1}{2x}-2$ 的最大值。
\hint{$x$ 与 $\dfrac{1}{2x}$ 均为负，先对 $-x$ 与 $-\dfrac{1}{2x}$ 用基本不等式，再把符号还原。与题 8 对照：只多一步转正。}
\end{problem}

\begin{problem}{题 10 （配凑定值）}
求下列函数的最小值：

(1) $y=x+\dfrac{1}{x-1}$（$x>1$）；\quad
(2) $y=\dfrac{x^2}{x-1}$（$x>1$）。
\hint{(1) 谁与 $\dfrac{1}{x-1}$ 相乘是定值？把 $x$ 写成 $(x-1)+1$；(2) 用 $x-1$ 改写分子：$x^2=(x-1)^2+2(x-1)+1$，逐项除开。}
\end{problem}

\begin{problem}{题 11 \starred（取倒数）}
求 $y=\dfrac{x-1}{x^2}$（$x>1$）的最大值。
\hint{它是题 10(2) 的倒数。$x>1$ 时 $y>0$，正数中较大者倒数较小，可直接引用 10(2) 的结论。}
\end{problem}

\begin{problem}{题 12 （和定积最大）}
求下列函数的最大值：

(1) $y=x(1-x)$（$0<x<1$）；\quad
(2) $y=x(1-2x)$（$0<x<\tfrac12$）。
\hint{(1) 两因子之和是多少？(2) 直接相加不是定值，先配系数：$y=\tfrac12\cdot 2x(1-2x)$，再看两因子之和。}
\end{problem}

\begin{problem}{题 13 （等号成立条件辨析）}
判断下列说法的真假，并说明理由：

A. $y=\sqrt{x^2+1}+\dfrac{4}{\sqrt{x^2+1}}$ 的最小值是 $4$；

B. 不等式 $\sqrt{x^2+1}+\dfrac{4}{\sqrt{x^2+1}}\ge 4$ 恒成立；

C. 对任意 $x>0$，$x^2+\dfrac{1}{x}\ge 2\sqrt{x}$ 恒成立；

D. 因为 $x=1$ 时 $x^2=\dfrac1x$（等号条件成立），所以 $y=x^2+\dfrac{1}{x}$（$x>0$）的最小值是 $y(1)=2$。
\hint{A、B 设 $t=\sqrt{x^2+1}$（注意 $t\ge 1$），等号要求 $t=2$，$x$ 取得到吗？C 与 D 形似而实异：C 只断言「$\ge$」，右边不必是定值；D 断言「最小值」，而 $2\sqrt{x^2\cdot\frac1x}=2\sqrt{x}$ 不是定值。取 $x=0.8$ 计算 $y$。}
\end{problem}

\clearpage
